% Acronym density -- title (left) and abstract (right).
% Data: ../../data/pgfplots/acronyms_title.csv, acronyms_abstract.csv

\begin{figure}[h]
\centering
\pgfplotsset{
  acro style/.style={
    height=5cm, width=0.47\linewidth,
    xlabel={Year}, ylabel={Acronyms per 100 words},
    xtick distance=4, x tick label style={rotate=45, anchor=east, font=\tiny,
                                          /pgf/number format/1000 sep={}},
    ytick align=inside, ymin=0,
    tick label style={font=\tiny}, label style={font=\small},
    enlarge x limits=0.02,
  }
}
% -- title density: legend collected here ----------------------------------
\begin{tikzpicture}
\begin{axis}[acro style,
  title={\small Titles~[\textdownarrow\ better]},
  legend to name=acrolegend,
  legend columns=5,
  legend style={draw=none, fill=none, font=\scriptsize,
                cells={anchor=west}, column sep=6pt},
]
  \addplot[color={rgb,255:red,44;green,160;blue,44}, dashdotted, very thick]
    table[x=year, y=neurips, col sep=comma] {../../data/pgfplots/acronyms_title.csv};
  \addlegendentry{NeurIPS}
  \addplot[color={rgb,255:red,148;green,103;blue,189}, solid, very thick]
    table[x=year, y=arxiv_ml, col sep=comma] {../../data/pgfplots/acronyms_title.csv};
  \addlegendentry{arXiv ML}
  \addplot[color={rgb,255:red,140;green,86;blue,75}, dotted, thick]
    table[x=year, y=arxiv_nonml, col sep=comma] {../../data/pgfplots/acronyms_title.csv};
  \addlegendentry{arXiv non-ML}
  \addplot[color={rgb,255:red,214;green,39;blue,40}, dashed, thick]
    table[x=year, y=pubmed, col sep=comma] {../../data/pgfplots/acronyms_title.csv};
  \addlegendentry{PubMed}
  \addplot[black, dashed, semithick, domain=1987:2025] {2.40};
  \addlegendentry{Barnett 2020 peak (2.40)}
  \draw[black, dashed, semithick] (axis cs:2022,0) -- (axis cs:2022,100);
\end{axis}
\end{tikzpicture}%
\hfill
% -- abstract density ------------------------------------------------------
\begin{tikzpicture}
\begin{axis}[acro style, title={\small Abstracts~[\textdownarrow\ better]}]
  \addplot[color={rgb,255:red,44;green,160;blue,44}, dashdotted, very thick]
    table[x=year, y=neurips, col sep=comma] {../../data/pgfplots/acronyms_abstract.csv};
  \addplot[color={rgb,255:red,148;green,103;blue,189}, solid, very thick]
    table[x=year, y=arxiv_ml, col sep=comma] {../../data/pgfplots/acronyms_abstract.csv};
  \addplot[color={rgb,255:red,140;green,86;blue,75}, dotted, thick]
    table[x=year, y=arxiv_nonml, col sep=comma] {../../data/pgfplots/acronyms_abstract.csv};
  \addplot[color={rgb,255:red,214;green,39;blue,40}, dashed, thick]
    table[x=year, y=pubmed, col sep=comma] {../../data/pgfplots/acronyms_abstract.csv};
  \addplot[black, dashed, semithick, domain=1987:2025] {2.40};
  \draw[black, dashed, semithick] (axis cs:2022,0) -- (axis cs:2022,100);
\end{axis}
\end{tikzpicture}

\vspace{2mm}
\pgfplotslegendfromname{acrolegend}

\caption{\textbf{Acronym density in titles (left) and abstracts
  (right), 1987 to 2025.}
  Lower values indicate fewer acronyms per word (\textdownarrow\
  better). Dashed black horizontal line marks the Barnett and
  Doubleday science-wide peak of 2.40 per 100 words (2019).
  NeurIPS and arXiv ML diverge from the non-ML control around
  2017 on titles and around 2020 on abstracts.}
\label{fig:acronyms}
\end{figure}
